\section{Semantics}
\label{sec:semantics}

The semantics of \texttt{MaybeUncertain<T>} are inherited almost entirely from
the underlying \texttt{Uncertain<T>} machinery. This section makes the
inheritance precise. We describe the computation graph that operators build,
the hypothesis test that the gating operator runs, and the way the two types
compose.

\subsection{The computation graph}
\label{sec:semantics:graph}

\texttt{Uncertain<T>} represents every value as a node in a directed acyclic
graph whose internal node enum is \texttt{UncertainNodeContent}. The variants
include leaf nodes (\texttt{Value}, \texttt{DistributionF64},
\texttt{DistributionBool}), arithmetic operators
(\texttt{ArithmeticOp}, \texttt{NegationOp}), logical operators
(\texttt{LogicalOp}), comparisons (\texttt{ComparisonOp}), conditionals
(\texttt{ConditionalOp}), and a small set of higher-kinded variants
(\texttt{PureOp}, \texttt{FmapOp}, \texttt{ApplyOp}, \texttt{BindOp}) that
the public API does not currently surface.

A \texttt{MaybeUncertain<T>} value is two such graphs, one rooted at
\texttt{is\_present} and one rooted at \texttt{value}. Operators construct
nodes in both graphs simultaneously. Binary arithmetic on
\texttt{MaybeUncertain<f64>} appends an \texttt{ArithmeticOp} node to the
value graph and a \texttt{LogicalOp::And} node to the presence graph; unary
negation appends a \texttt{NegationOp} to the value graph and leaves presence
unchanged. \Cref{fig:graph} sketches the combined structure.

\begin{figure}[h]
\centering
% TODO: TikZ sketch of two graphs (is_present, value) for c = a + b.
\fbox{\parbox{0.7\linewidth}{\centering [Figure to be drawn: Bayesian
network for $c = a + b$ where $a,b$ are \texttt{MaybeUncertain<f64>}. Two
panels: presence axis with an AND gate over the two Bernoulli leaves, value
axis with an addition node over the two value sub-graphs.]}}
\caption{Combined computation graph of $c = a + b$ over
\texttt{MaybeUncertain<f64>}. The presence axis builds a logical-AND graph;
the value axis builds the standard \texttt{Uncertain<T>} arithmetic graph.}
\label{fig:graph}
\end{figure}

\paragraph{Sampling the graph.} Each \texttt{Uncertain<T>} instance carries
a unique integer identifier assigned at construction. A sample is the result
of recursively evaluating the graph from the root, with leaf draws keyed by
$(\text{node id}, \text{sample index})$ in a thread-safe global cache so that
shared sub-graphs are evaluated once per index. Top-level \texttt{sample()}
draws a fresh random index from the thread-local generator on each call;
\texttt{sample\_with\_index} takes the index as an argument and produces a
deterministic draw for that pair.

The two-stage sampling of \texttt{MaybeUncertain<T>} described
in~\Cref{sec:type:sampling} sits on top of this machinery. The
\texttt{is\_present} graph is sampled with one index and, conditionally on a
true result, the \texttt{value} graph is sampled with a separately drawn
index. Independence between the two axes is therefore a property of the
implementation as written, not only of the type definition.

\subsection{The SPRT at the gate}
\label{sec:semantics:sprt}

\texttt{lift\_to\_uncertain} discharges its decision to
\texttt{Uncertain<bool>::to\_bool}, which in turn calls the implementation of
Wald's sequential probability ratio test~\cite{wald1945sprt,wald1948optimum}
shipped with the crate. The test as implemented compares a null and an
alternate hypothesis separated by an indifference region:
\begin{align*}
H_0 &: \Prob[\mathit{is\_present} = \text{true}] \le p_{\text{threshold}} - \varepsilon, \\
H_1 &: \Prob[\mathit{is\_present} = \text{true}] > p_{\text{threshold}} + \varepsilon.
\end{align*}
Here $p_{\text{threshold}}$ is the caller-supplied
\texttt{threshold\_prob\_some} and $\varepsilon$ is the caller-supplied
\texttt{epsilon}. The values $p_0 = p_{\text{threshold}} - \varepsilon$ and
$p_1 = p_{\text{threshold}} + \varepsilon$ are clamped to the open interval
$(0, 1)$ before use.

The test sets Type I and Type II error budgets to
$\alpha = \beta = 1 - \mathit{confidence\_level}$, where
\texttt{confidence\_level} is the third parameter to
\texttt{lift\_to\_uncertain}. Wald's two boundaries follow directly:
$$
A = \ln \frac{\beta}{1 - \alpha}, \qquad
B = \ln \frac{1 - \beta}{\alpha}.
$$
The test draws samples in batches of ten. After each batch it computes the
log-likelihood ratio for the $n$ samples drawn so far with $x$ successes,
\begin{equation*}
\Lambda_n = x \cdot \ln \frac{p_1}{p_0}
          + (n - x) \cdot \ln \frac{1 - p_1}{1 - p_0},
\end{equation*}
and compares to the boundaries. If $\Lambda_n \le A$, $H_0$ is accepted and
the gate returns false. If $\Lambda_n \ge B$, $H_1$ is accepted and the gate
returns true. Otherwise a new batch is drawn. If \texttt{max\_samples} is
reached without a decision, the implementation falls back to comparing the
empirical success rate $x/n$ to $p_{\text{threshold}}$ and returns the result
of that comparison.

Two consequences for callers matter in practice. The first is that termination
is not guaranteed in fewer than \texttt{max\_samples} batches; Wald's
guarantee is on expected sample size, and the fallback decision exists
precisely to bound the worst case. The second is that the indifference
region $[\,p_{\text{threshold}} - \varepsilon,\, p_{\text{threshold}} + \varepsilon\,]$
makes the test asymmetric in a useful way: presence probabilities inside the
indifference region produce slow decisions, and the empirical results
in~\Cref{sec:evaluation} confirm this regime is single-digit microseconds
slower than the easy-accept and easy-reject regimes.

\subsection{Composition with \texttt{Uncertain<T>}}
\label{sec:semantics:composition}

The two types are connected in both directions, although the connection is
not symmetric.

\paragraph{Injection.} Any \texttt{Uncertain<T>} embeds into
\texttt{MaybeUncertain<T>} via \texttt{from\_uncertain}. The injection sets
\texttt{is\_present} to a point distribution at true and copies the value
graph unchanged. Operations on the result behave exactly as on the original
\texttt{Uncertain<T>} value, except that the presence axis carries an AND of
\texttt{point(true)} leaves and therefore remains structurally true.

\paragraph{Projection.} The reverse direction is partial. A
\texttt{MaybeUncertain<T>} projects back to \texttt{Uncertain<T>} only when
\texttt{lift\_to\_uncertain} concludes that the SPRT supports the presence
hypothesis at the caller's chosen confidence; otherwise the projection
returns a \texttt{PresenceError}. The asymmetry is deliberate. An injection
adds information (``this value is definitely present''); a projection
discards information (the presence axis is collapsed to a single Boolean
verdict), and the discarding is only safe when statistical evidence
supports it.

\paragraph{Equational reading.} For any \texttt{Uncertain<T>} value $u$,
\texttt{lift\_to\_uncertain} on \texttt{from\_uncertain(u)} succeeds and
returns a clone of $u$, regardless of the confidence parameters. This holds
because the presence axis of \texttt{from\_uncertain(u)} is
\texttt{point(true)}, on which any SPRT against any threshold strictly less
than one trivially concludes acceptance after the first batch. The two
operations therefore compose as expected on the inhabited diagonal.

\paragraph{Predicates as observers.} \texttt{is\_some} and \texttt{is\_none}
expose the presence axis as a first-class \texttt{Uncertain<bool>}. Callers
can compose these with the existing \texttt{Uncertain<bool>} operators
(\texttt{\&}, \texttt{|}, \texttt{!}, \texttt{\^{}}) to build presence-aware
conditions over multiple \texttt{MaybeUncertain} values, and discharge those
conditions via the existing SPRT-backed methods such as
\texttt{probability\_exceeds} or \texttt{implicit\_conditional} on
\texttt{Uncertain<bool>}. No new machinery is needed.
